\documentclass[10pt,aspectratio=169]{ctexbeamer}
\usetheme{Boadilla}
\usecolortheme{seahorse}
\usepackage{amsmath,amssymb,booktabs}

\title{一个来源可追溯的数学报告}
\author{报告人}
\institute{单位}
\date{\today}

\begin{document}
\begin{frame}
  \titlepage
\end{frame}

\begin{frame}{目录}
  \tableofcontents
\end{frame}

\section{问题}
\begin{frame}{研究问题}
  \begin{block}{目标}
    简洁说明研究对象、数学问题和主要假设。
  \end{block}
  \begin{itemize}
    \item 已有结果是什么？
    \item 本报告的新内容是什么？
    \item 哪些假设不能省略？
  \end{itemize}
\end{frame}

\section{主要结果}
\begin{frame}{主要结果}
  \begin{theorem}
    在这里写出定理，并保留所有定义域、正则性和参数条件。
  \end{theorem}
  \begin{itemize}
    \item 来源：论文定理编号或笔记位置。
    \item 状态：已证明、猜想、实验观察或未来工作。
  \end{itemize}
\end{frame}

\section{证明思路}
\begin{frame}{证明或方法思路}
  \begin{enumerate}
    \item 先说明关键化简。
    \item 再给出核心构造或估计。
    \item 技术细节可放入 backup。
  \end{enumerate}
\end{frame}

\section{总结}
\begin{frame}{总结}
  \begin{itemize}
    \item 主要贡献。
    \item 主要限制。
    \item 下一步问题。
  \end{itemize}
\end{frame}
\end{document}
